\section{Addition with a log-general-type fiber}\label{generaltype:section}

We prove subadditivity when the fiber is of log general type.  The base
may be nonprojective: the projective positivity theorem will be applied
to a stable family over the image of its classifying map.

We prove \Cref{generaltype:addition}.

The proof constructs finitely many relative pluricanonical forms whose
ratios give a generically finite map on a general fiber.  Their poles
must be controlled at every vertical divisor, including divisors whose
image on the original base has codimension at least two.  Multiplying
these forms by base pluricanonical sections then retains both the fiber
coordinates and the base Iitaka coordinates.

We first construct a projective stable family over the image of a
classifying map.  Positivity is applied to that family.  Its sections
are pulled back, descended through a finite cover, and tested at
vertical valuations.  All degrees below clear the denominators of the
relevant divisors.

\subsection{Comparison with a stable family}

We use the stable-family convention of
\cite[Notation~3.7, Remark~3.8, and Definitions~3.9 and~3.11]
{KovacsPatakfalvi}. A stable pair has semi-log-canonical singularities
and an ample rational log canonical divisor. A stable family over a
normal base is flat, proper and surjective with connected stable
fibers; its boundary support avoids the generic points and the
codimension-one singular points of every fiber. Its relative log
canonical divisor is rational Cartier and relatively ample, so the
family is projective. Boundaries are pulled
back by divisorial restriction, and Cartier multiples of the relative
log canonical divisor commute with base change between normal bases.
Semi-log-canonicity
allows nonnormal fibers and tests log canonicity on the normalization
with the conductor boundary. The general fibers in our application
are additionally Kawamata log terminal (klt): they are normal and
all their log discrepancies are positive. This extra condition will
be checked before applying positivity.
A stable family has \emph{maximal variation} when its classifying map
to the moduli space of stable pairs is generically finite onto its image.

Fiberwise finite generation gives an embedding degree separately on
each fiber.  The next lemma selects one degree on an analytic base.
As throughout, \emph{very general} excludes a countable union of proper
closed analytic subsets.

\begin{lemma}\label{generaltype:selection}
Let \(Y\) be an irreducible compact complex space.  For each positive integer
\(j\), suppose that \(q_j\colon Q_j\to Y\) is a proper surjective map from an
irreducible compact complex space and that, over a dense analytic open
\(U_j\subset Y\), its fibers are irreducible of a fixed dimension \(r_j\).
Let
\[
  \phi_j\colon Q_j\dashrightarrow Z_j
\]
be a meromorphic map with irreducible projective image \(Z_j\).
Suppose that a dense open part of \(q_j^{-1}(U_j)\), called the genuine
frame locus, meets every fiber densely and that \(\phi_j\) is holomorphic
there.  Let \(G_j\subset Z_j\) be a countable union of algebraically
constructible subsets.

Let \(\Lambda_j\subset U_j\) be sets whose union contains a very general
subset of \(Y\), such that for \(y\in\Lambda_j\) all genuine frames above
\(y\) map into \(G_j\).  Then some \(\Lambda_j\) is not contained in a
countable union of proper closed analytic subsets of \(Y\), and for
that index one of the constructible subsets making up \(G_j\) is dense
in \(Z_j\).  In particular \(G_j\) contains a nonempty algebraic open
subset of \(Z_j\).
\end{lemma}

\begin{proof}
Some \(\Lambda_j\) is not contained in a countable union of proper
closed analytic subsets; otherwise their union would be so contained.
Resolve the graph of \(\phi_j\) and shrink \(U_j\), if necessary.
This discards only a proper analytic subset of \(Y\) and preserves
the stated properties of the general fibers.
Write \(G_j=\bigcup_iG_{ji}\).  If no \(G_{ji}\) is dense, the inverse
images
\[
  E_{ji}=\phi_j^{-1}(\overline{G_{ji}})
\]
are proper closed analytic subsets of the resolved \(Q_j\).
For \(y\in\Lambda_j\), the genuine frame fiber is covered by these
subsets.  Baire's theorem implies that one \(E_{ji}\) contains a
relatively open part of the fiber.  Irreducibility then implies
that \(E_{ji}\) contains the whole compact fiber.

The locus
\[
  A_{ji}=\{y\in U_j:
            \dim(E_{ji}\cap q_j^{-1}(y))\geq r_j\}
\]
is analytic, by the fiber-dimension theorem for proper maps.  It is
proper: otherwise \(E_{ji}\) would contain the general fibers and
hence all of \(Q_j\).  The same theorem before restriction to \(U_j\)
shows that \(A_{ji}\), together with \(Y\setminus U_j\), is contained
in a proper closed analytic subset of \(Y\).
This would place \(\Lambda_j\) in countably many proper closed
analytic subsets, a contradiction.  Finally, a dense constructible
subset of an irreducible algebraic variety contains a nonempty
open subset.
\end{proof}
\begin{lemma}[Analytic comparison with a stable family]
\label{generaltype:stable-comparison}
Let \(h\colon V\to Y\) be a surjective holomorphic map with connected
fibers, where \(V\) is a smooth compact Kähler manifold and \(Y\) is
a smooth compact manifold in class \(\mathcal C\).  Let \(A\) be an
effective rational SNC divisor on \(V\), horizontal over \(Y\), with
coefficients strictly less than one.  Suppose \(K_F+A_F\) is big on
very general fibers \(F\) of positive dimension.  There are
\begin{enumerate}[label=\textup{(\roman*)}]
\item a proper generically finite map \(\tau\colon Y'\to Y\), with \(Y'\)
smooth, compact, and in class \(\mathcal C\);
\item a surjective holomorphic map \(\rho\colon Y'\to R\), with \(R\)
smooth and projective;
\item a projective stable family
\(\pi\colon(\mathcal V,\mathcal A)\to R\) of maximal variation, with klt
general fiber;
\end{enumerate}
such that the main component of \(V\times_Y Y'\) is bimeromorphic
over \(Y'\) to \(\mathcal V\times_RY'\).  Over a dense analytic open,
this meromorphic comparison is the fiberwise log canonical model map
for \(K_F+A_F\), with the horizontal pushforward of \(A\) as boundary.
The space \(R\) may be a point.
\end{lemma}

\begin{proof}
The construction has three outputs.  Relative pluricanonical images
produce algebraic parameter data even though the base is analytic.
A classifying map then supplies a projective stable family after a
finite base cover.  Finally, an isomorphism cover matches that family
with the relative image; fiberwise canonical models determine the
horizontal discrepancy signs.

\emph{Relative images and frames.}
Set \(\mathcal L=\mathcal O_V(K_{V/Y}+A)\), understood as a rational
line bundle.  For the factorial multiples \(m\) clearing its denominator,
the sheaves
\[
  \mathcal E_m=h_*\mathcal L^{\otimes m}
\]
are coherent. First restrict to a dense analytic open where \(h\) and
the horizontal boundary strata are smooth, so the proper family and
its line bundles are flat over the base. Generic constancy of the
fiber section dimension and cohomology and base change then identify
the fibers of \(\mathcal E_m\) with the complete systems, separately
for each \(m\); see \cite[\S7, nos.~4--6, S\"atze~3--5]{GrauertDirectImages}
and its corrigendum \cite{GrauertDirectImagesCorr}.
After a modification of \(Y\),
the strict transform of \(\mathcal E_m\), modulo torsion, is locally free.
The evaluation map gives a relative meromorphic map into its projective
bundle.  Resolve this map and take its proper analytic image
\(\mathcal Z_m\).  We also take the proper images of the horizontal
boundary components and retain their divisorial cycle parts.
Analytic flattening makes these embedded image and cycle data flat on a
dense analytic open, with fixed Hilbert and cycle polynomials.
These uses of coherent direct images and analytic flattening are
independent of any minimal-model theorem; see \cite{HironakaFlattening}.

Write \(s_m=\rank\mathcal E_m\) on this open; indices with \(s_m=0\)
are omitted.  Its projective frame bundle has a natural compactification
\[
 \overline Q_m=
 \mathbb P\!\left(\Hom(\mathcal O_Y^{\,s_m},\mathcal E_m)\right).
\]
Here \(Y\) denotes the modification chosen for this degree; its map to
the original base remains proper and bimeromorphic.  The locus of
invertible matrices modulo scalars is a principal
\(\operatorname{PGL}(s_m)\)-bundle.
A genuine frame embeds the corresponding fiber of \(\mathcal Z_m\)
in a fixed \(\mathbb P^{s_m-1}\).  There is therefore a holomorphic
Hilbert--Chow parameter map on the genuine frame locus.
The Hilbert spaces here are the algebraic Hilbert schemes of
projective space, viewed analytically; equivalently one may use
their Douady realization \cite{Douady}.

To extend this map meromorphically, use the existing compact relative
image and the projective linear action in local trivializations.
The graph of that action extends over the projective matrix
space: inversion is given by the adjugate matrix.  Apply this construction
to the proper family \(\mathcal Z_m\) and its marked cycles, take proper
analytic images, and flatten by strict transform.  One obtains a proper flat family
of embedded subspaces over a modification of \(\overline Q_m\).
Its Hilbert--Chow map is holomorphic.  The local graph constructions
agree on the genuine frames and hence glue.  Thus the original
parameter map is meromorphic.  After resolving it, its image is
compact analytic in a fixed projective Hilbert--Chow space, and is
consequently algebraic.  Denote the irreducible image by \(Z_m\).
This construction uses a frame bundle, not a meromorphic
trivialization of \(\mathcal E_m\).

\emph{The good parameter locus.}
In \(Z_m\), call an embedded pair \((Z,D)\) good if it is klt and stable,
has the prescribed boundary weights, has a complete nondegenerate
embedding, and satisfies
\begin{equation}\label{generaltype:embedded-adjoint}
  \mathcal O_Z(1)\simeq
       \mathcal O_Z\bigl(m(K_Z+D)\bigr).
\end{equation}
One may also impose the generic Hilbert and marking data.
This is a countable union of algebraically constructible loci.
Here is the parameter-space justification.  Stratify the universal
embedded families by their Hilbert polynomials, Cartier indices,
and base-change loci.  On each such stratum,
\Cref{generaltype:embedded-adjoint} can be imposed by parameterizing
the map
\[
  \omega_Z^{[m]}\longrightarrow\mathcal O_Z(1)
\]
with its prescribed divisor.  These parameter spaces are of finite
type.  The stable-pair and klt conditions are constructible there,
and their images are constructible by Chevalley's theorem.
This is the embedded construction used for stable log-variety
moduli in \cite[Section~6.A, especially Propositions~6.8 and~6.11
and Lemma~6.12]{KovacsPatakfalvi}.  In particular no assertion
about analytic minimal models is hidden in the good locus.

Every very general fiber is projective: bigness makes it Moishezon,
and it is Kähler. We use Moishezon's projectivity criterion, which
is also the smooth case of \cite[Theorem~6]{NamikawaProjectivity}.
The projective klt finite-generation theorem
\cite[Theorem~1.1 and its consequences]{BCHM} therefore gives its log
canonical model.  Every sufficiently large divisible pluricanonical
degree embeds that model and satisfies
\Cref{generaltype:embedded-adjoint}.  Exclude at once the countably
many generic base-change loci, relative-rank loci, and generic-degree
loci for the maps constructed above.  Let \(\Lambda_m\) consist of
the remaining parameters for which the degree-\(m\) map actually
constructs the fiber's log canonical model.
Their union contains a very general subset of \(Y\), and all genuine
frames over a point of \(\Lambda_m\) give good data.
Apply \Cref{generaltype:selection}.  We obtain one fixed \(m\) for which
\(\Lambda_m\) is not contained in a countable union of proper analytic
subsets and good data hold on an algebraic open of \(Z_m\).

For this degree the relative image map is generically birational.
Indeed its generic finite degree agrees with its fiberwise degree
outside the generic-degree exceptional set already excluded.
Choose a point of \(\Lambda_m\) there; that degree is one.
Retain \(\Lambda_m\): its actual canonical models will determine the
horizontal discrepancy signs at the end of the construction.

\emph{The invariant classifying map.}
Fix the dimension, the volume read from
\Cref{generaltype:embedded-adjoint}, and a finite coefficient set
containing the coefficients of \(A\) and closed under sums at most one.
The good algebraic parameter locus has a classifying map to the
projective coarse space \(M\) of stable pairs with these data.
If an additional parameterization of the map in
\Cref{generaltype:embedded-adjoint} was used, its moduli graph descends:
the embedded pair determines a unique isomorphism class, so the
closure of that graph is generically a singleton over the good
parameter locus.  Restricting a dense algebraic open therefore gives
the same rational classifying map on \(Z_m\).
It is invariant under change of projective frame.

Compose with the meromorphic frame parameter map.
The resulting map \(\overline Q_m\dashrightarrow M\) is constant
on general frame fibers.  Its graph has a proper analytic image
in \(Y\times M\).  This image is generically a singleton over \(Y\),
so its projection to \(Y\) is bimeromorphic.  It is the graph of
a meromorphic map
\[
  \mu\colon Y\dashrightarrow M.
\]
Let \(M_0\) be the closure of its image; it is projective.

\emph{A stable parameter cover and the matching cover.}
Stable moduli supplies a finite parameter cover
\(S\to M\) from a normal scheme carrying a stable family
\cite[Corollary~6.19]{KovacsPatakfalvi}.
The coarse space \(M\) is projective
\cite[Corollary~7.3]{KovacsPatakfalvi}, so its finite cover \(S\)
is projective as well. The moduli construction for the specified
coefficient set is given in
\cite[Definition~6.2, Proposition~6.4, and
Proposition~6.14]{KovacsPatakfalvi}.
Resolve \(\mu\), and normalize a component of the pullback of
\(S\to M\) dominating \(Y\).  Let \(R_0\) be its image in \(S\).
The restricted stable family over \(R_0\) has maximal variation,
because \(R_0\to M_0\) is finite.

At this point the projective family exists over \(R_0\).
Equality of coarse moduli points identifies its general fibers with
those of the original analytic image, but does not identify the
families. We construct a further generically finite cover carrying
such an identification.

On a dense analytic open of the comparison base, embed the original
flat image by the complete system of
\(\mathcal O_{\mathcal Z_m}(k)\), and the pulled-back stable family
by its \(km\)-adjoint system, for a common sufficiently large and
divisible \(k\). The two sheaves of sections are coherent direct
images of compact families; generic cohomology and base change make
them locally free after shrinking. On a common compact modification
of the base, make their torsion-free strict transforms locally free,
resolve the evaluation maps, and flatten the embedded image families
and their weighted boundary cycles. This gives holomorphic relative
Hilbert--Chow data. The modification concerns the base itself, not a
space of choices of frames.

On every general good fiber, \Cref{generaltype:embedded-adjoint}
identifies the two polarizations. Thus the projective linear maps
identifying the two embedded varieties and their weighted boundary
cycles form a nonempty Isom set. This set is finite, since stable
pairs have finite automorphism groups
\cite[Propositions~6.5 and~6.8]{KovacsPatakfalvi}; completeness and
nondegeneracy of the embeddings exclude additional ambient
stabilizers. At this stage we use only the embedded analytic image
and its polarization. Its \(\Q\)-Gorenstein family structure will
be transported from the stable family by the identification.

To compactify these Isom sets, work in local trivializations of the
two section bundles. In the product of their projective parameter
spaces with the projective matrix space, take the algebraic closure
of the incidence relation for invertible matrices identifying the
embedded varieties and weighted boundary cycles, after clearing the
fixed coefficient denominators. Pull back this closure by the two
Hilbert--Chow maps. Equivariance under changes of both frames glues
the local incidence spaces into a proper analytic space in the
projective bundle of homomorphisms between the section bundles; no
meromorphic trivialization is used. Retain its compact irreducible
components that dominate the base and meet the invertible locus,
and normalize them. The invertible locus has the finite Isom fibers
just described, so the resulting cover is proper and generically
finite. Its tautological isomorphism identifies the generic families,
including the stable family's \(\Q\)-Gorenstein structure; its graph
therefore extends meromorphically over the compact base.

Finally resolve \(R_0\) projectively and resolve the graph of the
map from the comparison cover to it.  Denote the resulting spaces
by \(R\) and \(Y'\).  Stable families and their divisible relative
adjoint line bundles commute with base change, with boundary understood
by divisorial restriction \cite[Remark~3.8]{KovacsPatakfalvi}.
Thus the pullback family over \(R\) remains stable; it has klt general
fibers and maximal variation.  All covers and modifications used above stay
in class \(\mathcal C\).

On a common resolution we can now compare the source adjoint
divisor with the pullback of the stable relative adjoint divisor,
which is rational Cartier.  There are only finitely many horizontal
prime coefficients in their difference.  On fibers above
\(\Lambda_m\), away from the finitely many generic comparison
exceptions, they are the log canonical model discrepancy coefficients.
They are therefore nonnegative and supported on divisors exceptional
over the stable model.  The same statements hold horizontally in
the family.  Remove the images of the remaining vertical discrepancy
divisors.  On the resulting dense analytic open this is the log
canonical model comparison, with its divisorial pushforward boundary.
This proves the asserted meromorphic comparison without relative
analytic finite generation.  If \(M_0\) is a point, the construction
uses the fixed stable pair and \(R\) is a point.
\end{proof}

\subsection{Relative forms and vertical valuations}

For a meromorphic section \(\sigma\) of
\(\mathcal O_V(bK_{V/Y})\), its horizontal pole bound \(bA\)
means that \(\ord_E(\sigma)\geq-bA_E\) at every horizontal prime,
using a local nonvanishing base volume to interpret relative forms.
At a vertical prime over a divisor \(D\subset Y\), write
\(\omega_D\) for a local nonvanishing logarithmic top form on the
base with a simple pole along \(D\), and no other pole at its generic
point.  We say that \(\sigma\) has the relative logarithmic bound at
this prime if
\begin{equation}\label{generaltype:relative-log-bound}
  \ord_E\bigl(\sigma\wedge h^*\omega_D^{\,b}\bigr)\geq-b.
\end{equation}
The notation \(\wedge\) denotes the relative-times-base identification
of the corresponding pluricanonical line bundles.

\begin{lemma}\label{generaltype:stable-orders}
In the situation of \Cref{generaltype:stable-comparison}, pull back a
holomorphic section of a divisible multiple of
\(K_{\mathcal V/R}+\mathcal A\), and compare it with relative pluriforms
on the main component of
\(V\times_Y Y'\).
It satisfies the horizontal bound given by the pullback of \(A\).
Moreover, on a common resolution it satisfies
\Cref{generaltype:relative-log-bound} at every prime lying over a
divisor of any smooth modification of the comparison base.
\end{lemma}

\begin{proof}
On general fibers, pullback from the log canonical model lies in the
original log pluricanonical system.  The effective canonical-model
discrepancy on a common resolution therefore proves the horizontal bound.

For the vertical assertion, first fix a smooth modification of the
comparison base and a prime divisor on it.  The stable family and its
divisible relative adjoint line bundles commute with this base change.
Work transversely to a general point of that divisor.  The resulting
one-parameter family is still stable.  This order of operations is
what permits the same bound at divisors exposed only on a higher base
model.  Its total pair with the reduced
central fiber added is log canonical.  More explicitly, choose the
disk with klt fibers away from its origin.  Its total space is normal:
flatness over the smooth disk and the \(S_2\) property of the fibers
give \(S_2\).  A codimension-one point on the central fiber is a
generic point of one of its reduced components.  Its local ring is
one-dimensional, and its quotient by the disk parameter is a field;
its maximal ideal is consequently generated by that parameter, so
this local ring is regular.  Away from the central fiber the total
space is normal, because its fibers are klt.  Serre's criterion now
gives normality.
Compatibility of the relative adjoint divisor with the stable fiber
identifies the adjunction different on the normalization of the central
fiber with its stable boundary and conductor.  The resulting pair is
log canonical because the central fiber is stable; see
\cite[Lemma~2.10 and Corollary~2.11]{PatakfalviFibered} for this
normalization-and-conductor identity.
The analytic Cartier-divisor inversion theorem
\cite[Theorem~1.1]{FujinoAnalyticAdjunction} consequently applies.
The argument also applies after any finite ramified disk change.

Let \(t\) be the disk parameter and let \(\xi\) denote the
pulled-back relative section in degree \(b\).  Then
\(\xi\wedge(dt/t)^b\) is a section of the total log canonical
multiple with the central fiber added.  On a log resolution the
log canonical discrepancy bound says that this form has pole
order at most \(b\) at every vertical prime.  A further common
resolution with the original family preserves this bound.
This is \Cref{generaltype:relative-log-bound}.  Smooth transverse
parameters do not change any of these orders.
\end{proof}

\begin{lemma}\label{generaltype:relative-system}
Under the hypotheses of \Cref{generaltype:stable-comparison}, there
is a positive integer \(b_0\) and a finite-dimensional space
\[
  \mathcal S\subset
  H^0_{\mathrm{mer}}\bigl(V,\mathcal O_V(b_0K_{V/Y})\bigr)
\]
whose associated meromorphic map has rank \(\dim F\) on general
\(h\)-fibers, with the following properties:
\begin{enumerate}[label=\textup{(\roman*)}]
\item every member has horizontal pole bound \(b_0A\);
\item every member satisfies the relative logarithmic bound
at primes over base divisors;
\item the same bound may be tested after any modification of
the base, at a divisor exposed on that modification.
\end{enumerate}
The same assertions hold for the systems of products of members
of \(\mathcal S\).
\end{lemma}

\begin{proof}
Use \Cref{generaltype:stable-comparison}.  If \(\dim R>0\), the
stable-family bigness theorem
\cite[Theorem~8.1 and Corollary~8.3]{KovacsPatakfalvi} applies:
the base is normal projective, the family is stable and has
maximal variation, and its general fiber is klt.
Thus \(K_{\mathcal V/R}+\mathcal A\) is big.
Choose finitely many sections of one sufficiently divisible
multiple whose map is generically finite.  Its image has dimension
\(\dim F+\dim R\), and its restriction to a general \(\pi\)-fiber
has rank \(\dim F\).  If \(R\) is a point, choose such
sections of the ample log canonical divisor of the fixed stable
pair instead.  Pullback and the meromorphic comparison preserve these
image dimensions and give the pole bounds of
\Cref{generaltype:stable-orders}.

Replace the comparison cover by a Galois cover and resolve the
diagrams equivariantly.  To do this analytically, first take its Stein
factorization over \(Y\); its finite factor has the same cover over a
dense open set.  Over the finite étale locus, the Galois closure is a
component of a sufficiently large fiber product of this finite factor.
Take the corresponding compact component, normalize, and then resolve
the graph of its map to the comparison cover.  Equivariant resolution
gives the required diagram; write \(G\) for its finite group.
This construction uses finite analytic maps and does not require the
meromorphic function field of \(Y\) to have transcendence degree
\(\dim Y\).

Let \(\mathcal B\) be the graded \(\C\)-algebra generated by the
chosen finite set of relative sections and all its \(G\)-conjugates.
Every homogeneous element has the required horizontal and vertical
linear pole bounds: products add the orders and sums cannot
decrease the lower bounds.  The algebra is integral and finite over
\(\mathcal B^G\).  Indeed each homogeneous generator satisfies its
orbit polynomial with invariant coefficients, and invariants of
a finite group acting on a finitely generated algebra are finitely
generated.  The corresponding map
\(\operatorname{Proj}\mathcal B\to\operatorname{Proj}\mathcal B^G\)
is finite.  Before taking invariants, the system contains the sections
pulled back from \(\mathcal V\), whose restriction to a general source
fiber has image dimension \(\dim F\).  A finite map preserves the
dimension of the image of that fiber, even though \(G\) may permute
the fibers above one point of \(Y\).  Choose a Veronese degree for
\(\mathcal B^G\) generated by finitely many elements of that single
degree.  Their ratios therefore still have rank \(\dim F\) on general
fibers.  Denote their pluricanonical degree by \(b_0\).
The same finite quotient also preserves the total image dimension,
which is at least \(\dim F+\dim R\).

Use the natural identification of relative canonical bundles on the
common smooth locus for every base change or modification.  On the
dense étale locus, invariant relative
tensors therefore descend to \(V\).  They extend meromorphically
across the remaining locus by finite analytic descent.  The
horizontal bounds descend because the comparison cover is
generically étale along each horizontal divisor.

For completeness, the vertical order calculation is exact.
Let \(E'\) lie over \(E\), with ramification index \(e\).
For a degree-\(b_0\) top pluriform \(\alpha\),
\begin{equation}\label{generaltype:ramified-order}
  \ord_{E'}(\alpha_{\mathrm{pullback}})+b_0
       =e\bigl(\ord_E(\alpha)+b_0\bigr).
\end{equation}
A logarithmic top form on the base pulls back to a nonvanishing
logarithmic top form at a prime dominating its divisor.
Apply \Cref{generaltype:ramified-order} to the relative form
wedged with this base logarithmic form.
The bound upstairs is therefore equivalent to
\Cref{generaltype:relative-log-bound} downstairs.
For (iii), pull the chosen Galois comparison cover and its original
sections to a common refinement over the modified base.
Lemma~\ref{generaltype:stable-orders} applies to these same pullbacks
at primes above the newly exposed base divisor; its bounds pass to
their conjugates, products, and invariant combinations.  The
ramification calculation therefore gives the bound for the pullback
of each fixed descended member of $\mathcal S$.
Products add the order inequalities, and taking a full product system
is a Veronese operation on its image, so products preserve the relative
rank as well.
\end{proof}

\begin{lemma}\label{generaltype:expose-valuation}
Let \(h\colon V\to Y\) be a proper surjective holomorphic map
of irreducible compact complex spaces, with \(Y\) smooth.
If a prime divisor \(E\) of a smooth \(V\) does not dominate \(Y\),
there is a modification of the base and a dominating main-component
model on which the strict transform of \(E\) maps onto a prime
divisor of the modified base.
\end{lemma}

\begin{proof}
Apply analytic flattening to \(h\), by strict transform, and resolve
the modified base.  Pullback preserves the flatness of the intermediate
family, whose fibers have the generic relative dimension \(d\).
This intermediate source is proper and bimeromorphic over \(V\).
On a resolution dominating it, the strict transform of \(E\) remains
a divisor, because a proper modification of the smooth \(V\) is an
isomorphism at the generic point of \(E\).  Its image in the flat
intermediate source has dimension \(\dim V-1\): it still dominates
\(E\).  The fiber-dimension bound consequently gives an image on the
new base of dimension at least
\(\dim V-1-d=\dim Y-1\).
That image is proper, since its projection to \(Y\) is contained in
\(h(E)\), and properness makes it a closed analytic set.  It is
therefore a prime divisor.  Further source resolutions retain the
same generic image and do not change the valuation of \(E\).
\end{proof}

\subsection{Proof of the addition proposition}

\begin{proof}[Proof of \Cref{generaltype:addition}]
We may suppose \(\kappa(Y,K_Y+C)\geq0\).
Choose a Kähler modification \(\nu\colon\widetilde V\to V\),
resolve the boundary,
and put
\[
  \widetilde T=\nu_*^{-1}T+\operatorname{Exc}(\nu)_{\mathrm{red}}.
\]
By \Cref{setup:log-modification,setup:modification-compatibility},
pullback and descent identify the divisible log pluricanonical rings
on the total spaces and on their very general fibers.  In particular,
fiberwise bigness is preserved, and
\(B(h\circ\nu,\widetilde T)=B(h,T)\).
Every prime over base codimension at least two still has coefficient
one: it is either a strict transform of such a prime or is newly
exceptional.  Write \((V,T)\) for this Kähler model; its sections
descend to the original pair.

If \(\dim F=0\), the map is bimeromorphic.  Pull back sections
of \(K_Y+C\).  At a prime dominating a base divisor the
multiplicity is one and the condition \(C\leq B(h,T)\) gives
the required pole bound.  At an exceptional prime the
log canonical discrepancy bound for \((Y,C)\), together with
\(T_E=1\), gives the same conclusion.  Thus pullback preserves
all the base systems and proves the assertion.  If \(Y\) is a
point, the assertion is simply the assumed bigness.  We may
therefore assume both dimensions are positive.

Write \(T^{\mathrm h}\) for the horizontal part of \(T\).
Choose a rational \(\epsilon>0\) sufficiently small and set
\[
  A=(1-\epsilon)T^{\mathrm h}.
\]
Its coefficients are strictly less than one.
To make \(K_F+A_F\) big uniformly on very general fibers, consider
\(\epsilon=1/j\) and all denominator-clearing degrees.  Choose a fiber
outside their countably many generic base-change and relative-rank
exceptional sets.
Openness of the big cone on that projective fiber gives one
choice of \(\epsilon\) and one degree with full image dimension.
The generic-rank choice makes the same degree have full rank
on general fibers.  Notice that no vertical coefficient has
been lowered, and \(C\) has not been changed.

Apply \Cref{generaltype:relative-system} to \(A\), and write
\(\sigma\) for a member of its system in degree \(b\), including
any product degree.  Let
\[
  \eta\in H^0\bigl(Y,\mathcal O_Y(b(K_Y+C))\bigr).
\]
The relative-times-base product
\(\sigma\wedge h^*\eta\) is a meromorphic section of
\(\mathcal O_V(bK_V)\).  We check that its poles are bounded by
\(bT\) at every prime.

There is nothing further to check at horizontal primes, since
\(A\leq T^{\mathrm h}\).  Suppose next that \(E\) maps onto a
prime divisor \(D\subset Y\), and put
\[
  a_E=\ord_E(h^*D).
\]
Relative to a local nonvanishing logarithmic base volume
\(\omega_D\), the section \(\eta\) vanishes to order at least
\(b(1-C_D)\).  The relative logarithmic bound gives
\begin{equation}\label{generaltype:orbifold-order}
  \ord_E(\sigma\wedge h^*\eta)
        \geq -b+b\,a_E(1-C_D).
\end{equation}
The inf-multiplicity convention is precisely
\[
 B(h,T)_D
   =1-\max_{E\mapsto D}\frac{1-T_E}{a_E}.
\]
Consequently \(C_D\leq B(h,T)_D\) implies
\[
 a_E(1-C_D)\geq1-T_E,
\]
and \Cref{generaltype:orbifold-order} is at least \(-bT_E\).
This uses no integrality condition on the boundary multiplicities.

Finally, suppose \(\operatorname{codim}_Y h(E)\geq2\).
Expose a base divisor \(D'\) below the strict transform of \(E\)
by \Cref{generaltype:expose-valuation}, and resolve the diagrams.
The pullback of \(\eta\) has at most a logarithmic pole along
\(D'\): this is the divisorial discrepancy inequality for the
log canonical smooth pair \((Y,C)\).
Property (iii) of \Cref{generaltype:relative-system} therefore
gives order at least \(-b\) at the strict transform of \(E\).
The order is unchanged, since source modifications are isomorphisms
at the generic point of \(E\).
The assumption \(T_E=1\) supplies exactly this allowance.
We have proved, at all prime divisors, that
\begin{equation}\label{generaltype:product-sections}
 \sigma\wedge h^*\eta
       \in H^0\bigl(V,\mathcal O_V(b(K_V+T))\bigr).
\end{equation}
On a smooth space these divisorial bounds imply holomorphic
extension of the corresponding log pluricanonical section.

Choose a degree in which the complete base system has image
dimension \(\kappa(Y,K_Y+C)\), and pass to a common multiple
with \(b_0\) in \Cref{generaltype:relative-system}.
The product system of relative sections still has image dimension
\(d=\dim F\) on a general fiber.  The corresponding product of the
chosen base system has image dimension \(\kappa(Y,K_Y+C)\), since a
Veronese embedding preserves its image.  In this common degree write
the nonzero relative generators as \(\sigma_i\) and the base
generators as \(\eta_j\).
The total system in \Cref{generaltype:product-sections} contains
all \(\sigma_i\wedge h^*\eta_j\).
On the open set where \(\sigma_0\) and \(\eta_0\) are nonzero,
ratios in this system include
\[
  \frac{\sigma_i}{\sigma_0},
  \qquad
  h^*\!\left(\frac{\eta_j}{\eta_0}\right).
\]
Its image therefore maps onto the base Iitaka image, and over
a general point of that image it contains the \(d\)-dimensional
image of a general \(h\)-fiber.
The dimension-of-fibers inequality gives image dimension at
least \(d+\kappa(Y,K_Y+C)\).  This proves
\Cref{generaltype:inequality} on the Kähler model.
The initial log pluricanonical descent then proves it on the
given smooth source.
\end{proof}
